% SEGMENT OLP-0604-S01
% Part: methods
% Chapter: proofs
% Section: starting-proofs

\documentclass[../../../include/open-logic-section]{subfiles}

\begin{document}

\olfileid{mth}{prf}{str}

\olsection{நிறுவலைத் தொடங்குதல்}

எங்கிருந்து தொடங்குவது?

நீங்கள் ஒன்றை நிறுவ வேண்டும் என்று தரப்பட்டுள்ளது. எனவே
நிறுவலில் கடைசியாக வர வேண்டியது அதுதான். தொடக்கத்தில் அதை
நிறுவப்போவதாக \emph{அறிவிக்கலாம்}; ஆனால் அதையே எடுகோளாகப்
பயன்படுத்தக் கூடாது. புதிய தாளின் கீழ்ப்பகுதியில் நிறுவ
வேண்டியதைக் குறித்துக்கொள்ளுங்கள். அப்போது இலக்கு உங்கள்
கவனத்திலிருந்து மறையாது.

% SEGMENT OLP-0604-S02
அடுத்து, பயன்படுத்தக்கூடிய சில எடுகோள்கள் இருக்கலாம்.
அடுத்த பிரிவில் நீங்கள் செய்யும் நிறுவலின் \emph{வகையைப்}
பற்றிப் பேசும்போது இது இன்னும் தெளிவாகும். அவற்றைத் தாளின்
மேற்பகுதியில் எழுதுங்கள்; அவை எடுகோள்கள் என்பதைத் தெளிவாகக்
குறியுங்கள். எடுத்துக்காட்டாக, $p$ ஐ எடுகோளாகக் கொண்டால்,
``$p$ எனக் கொள்வோம்'' அல்லது ``$p$ என்க'' என்று எழுதுங்கள்.
இறுதியாக, கேள்வியில் நீங்கள் அறிந்திருக்க வேண்டிய சில
வரையறைகள் இருக்கலாம். குறிப்பிட்ட வரையறையைப் பயன்படுத்தச்
சொல்லியிருக்கலாம்; அல்லது எடுகோள்களிலும் நீங்கள் அடையவேண்டிய
முடிவிலும் பல வரையறைகள் இருக்கலாம். \emph{அவற்றை எழுதிவைத்து,
அவை என்ன பொருள் தருகின்றன என்பதை உறுதிப்படுத்திக்கொள்ளுங்கள்.}

% SEGMENT OLP-0604-S03
நிறுவலை எப்படி அமைப்பது என்பது கேள்வியின் வடிவத்தையும்
சார்ந்திருக்கும். வாக்கியத்தின் வகைக்கேற்ப நிறுவலை எப்படி
அமைப்பது என்பதை அடுத்த பிரிவு விளக்குகிறது.

\end{document}
